\section{Preprocessing and Building Training Examples}
\label{sec:timeline-and-milestones}

\subsection{Tokenization Pipeline}

The process follows a three-stage pipeline for each sample:

\begin{enumerate}
    \item \textbf{Prompt Rendering:} The math problem is formatted into a chat prompt (e.g., using \texttt{<|im\_start|>user} tags).
    \item \textbf{Answer Formatting:} The teacher's reasoning trace and final answer are combined into a single target string.
    \item \textbf{Concatenation:} The prompt and answer are joined into a single token sequence ending with an end-of-sequence (\texttt{<|im\_end|>}) token.
\end{enumerate}

\subsection{Filtering and Sequence Management}

Reasoning traces can be excessively long, which increases computational requirements.

\begin{itemize}
    \item \textbf{Average Length:} The initial dataset averaged 2,946 tokens per response.
    \item \textbf{Outliers:} Some traces reached up to 42,005 tokens.
    \item \textbf{Filtering Strategy:} To keep costs reasonable and fit within hardware constraints, the dataset was filtered to a maximum length of 2,048 tokens. This removed 5,305 of the original 12,000 examples, resulting in a training set of 6,695 high-quality examples.
\end{itemize}